\section{The Model in Brief}
\label{sec:model}

We recall only what is needed to read the results. The full derivation is in Pollard (2026).

\begin{definition}[Vibe, tone, note, mesh]
A \emph{vibe} is the one kind of entity, a state in relation. Each vibe carries a finite \emph{tone}, with the working alphabet ternary, $t \in \{-1, 0, +1\}$. There is no richer tone on a single vibe. The spinor of matter (Section~\ref{sec:matter}) is a composite pattern of ternary vibes, not a richer primitive alphabet, with its complex phase emergent rather than fundamental. The \emph{note} is the only primitive relation, $v$ noting $w$ meaning $v$ is present to $w$. The \emph{vibe mesh} is the whole population of vibes, their tones, and the note relation among them. There is no space or time outside the mesh.
\end{definition}

The tone is not a separate thing inside a vibe. A vibe is a unit of experience, and its tone is simply the quality of that experience, the vibe seen as a felt state, not a component bolted onto it. We say a vibe carries a tone only as a manner of speaking, the way one names the colour of a surface without supposing the colour is an object apart from the surface. The ternary values are three readings of that one quality, not three contents stored in a slot.

The note is not a separate kind of object. By monism there is only the vibe, so a note is nothing inserted between vibes. It is simply a vibe's experiencing of other vibes, what it notes, the relation a vibe already bears by being present to others. We give that relation a name, the note, only to make the structure easier to model, discuss, and reason about, the way one names an edge in a graph without supposing the edge is a thing apart from the nodes it joins.

It helps to hold two complementary views of the same fact. A \emph{note} and a \emph{link} are the same vibe-to-vibe relationship seen from two sides. The note is that relationship from the inside, a vibe experiencing what it notes, the first-person reading. The link is the same relationship from the outside, an edge in the web, the structural reading used when we draw the mesh as a graph and run operators on it. They are similar but not identical lenses on one underlying interconnectedness, and this paper uses whichever is clearer for the structure at hand, the note when the point is experiential or causal, the link when the point is geometric or computational. Neither lens adds a new substance.

When a relation needs to carry its own state, that state is itself held by a vibe, never by a separate link. So everywhere below that we speak of notes, links, or the web, the underlying ontology is one of vibes experiencing vibes, and the relational vocabulary is a modelling convenience laid over that single substance.

The mesh grows and updates by a local \emph{rule}, one \emph{beat} at a time, each vibe feeling only its neighbors. Four committed choices shape it, each with named alternatives held open:

\begin{itemize}
\item \textbf{Discreteness.} The base is finite. Real numbers appear only as emergent large-scale measurements, never in the substrate.
\item \textbf{Hyperbolic geometry.} The mesh branches outward exponentially rather than as a flat grid.
\item \textbf{Eternal expansion.} The mesh keeps growing without bound.
\item \textbf{No hidden state.} Everything is in the tones and notes. There is no separate bookkeeping behind them.
\end{itemize}

The causal order of the mesh, the relation of one vibe being in another's past, is a locally finite partial order, which is exactly a \emph{causal set} (Bombelli et al., 1987). This identification is what lets the geometry, dynamics, and matter of the mesh be studied with the established discrete-spacetime toolkit, and it is the bridge between Vibe Theory's ontology and the rest of this paper.

The guiding picture for the results is a ladder. At the bottom is the mesh of tones and notes. Above it, as large-scale patterns, are space, time, matter, force, gravity, and the correlations we call quantum. The paper climbs this ladder one rung at a time, and the final section shows how the rungs connect into a single structure.
